%!TEX root = ../../main.tex
\section{방향, 평균, 규모는 다른 객체이다}
\label{sec:val-triad}

교전을 $\ppre$ 구간으로 나누어 보면 부호 있는 평균 변화 $\E[\dV]$, 절대 규모 $\E|\dV|$, 양의 방향 비율 $\Pr(\dV > 0)$이 함께 움직일 필요가 없다. 균형에 가까운 $\ppre$에서는 $\E[\dV]$가 0에 가까우면서 $\E|\dV|$는 가장 크고, 높은 $\ppre$(예: $[0.9, 1.0]$)에서는 $\Pr(\dV > 0) \approx \triadHighP$이면서 $\E[\dV]$는 약간 음이다. 이는 변화 방향 예측 문헌 \cite{ChristoffersenDiebold2006}이 말하는 부호와 평균의 분리를 이 자료에서 그대로 보여 준다. 따라서 이진 $\SVI$를 주 목표로 유지하되, 큰 규모를 방향이 교전 전 상태로 예측 가능하다는 증거로 다루지 않는다. 이 삼중 분해(triad)는 예측 가능성을 증명하지 않으며 측정 대상이 무엇인지 보여 줄 뿐이다.
